\documentclass[11pt,letterpaper]{article}
\usepackage{styles/michaillat-paper}
\usepackage{styles/working-paper}
\usepackage{styles/yax-common}
\geometry{margin=1in}
\hypersetup{pdftitle={Revision Diagnosis},pdfauthor={Lily Luo}}
\setlength{\parindent}{0pt}
\setlength{\parskip}{0.6em}
\begin{document}
\begin{center}
{\Large\bfseries Revision Diagnosis}\par
\vspace{0.3em}
Evidence-led framing decision, completed work, and unresolved limits\par
September 2026
\end{center}

\section*{Framing decision}

The defensible paper is a \textbf{comparison-centered empirical audit}, not an AI-displacement design and not a generic ``measurement choices matter'' essay.  Its question is: what occupational comparisons support the negative association between GPT-4 task exposure and young-relative employment stock in the public CPS?

The answer is specific.  The reconstructed pooled Q5--Q1 comparison is $-0.1321$ log point, with 95-percent occupation wild-score interval $[-0.2206,-0.0437]$.  Giving every broad occupation family its own young-relative calendar-month path yields $-0.0217$, with interval $[-0.1607,0.1173]$.  Their paired difference is 0.1104, with interval $[0.0107,0.2102]$.  The movement is detected under the declared occupation-shock procedure, but it compares different conditional estimands and is not a causal share.

Direct support is the binding economic qualification and the paper's organizing contribution.  Only four of 22 broad families contain both exposure tails; their 29 tail occupations account for 5.03 percent of preperiod stock, and the direct-tail comparison has an effective information count of 6.2 occupations.  Other families connect the coefficient through intermediate quintiles and common-coefficient restrictions.  The article should therefore say that the pooled pattern is principally a comparison across broad occupational structure, not that ``composition explains'' a known fraction of an AI effect.

\section*{Evidence that determines the interpretation}

\begin{enumerate}[leftmargin=1.8em]
\item \textbf{Corrected baseline.}  Five wrongly selected March ASEC samples are replaced by Basic Monthly files; October 2025 is absent; December 2022 is the transition month.  All 71-month preperiod weights, normalizations, cuts, ties, and memberships are rebuilt.  Nine occupations change quintile, but reconstructed-minus-historical treatment moves the coefficient only 0.00245, with interval spanning zero.
\item \textbf{Support geometry, not attenuation alone, carries the contribution.}  The pooled-to-SOC2-month movement is 0.1104 and statistically distinguishable from zero on common draws, while the conditioned coefficient and direct-tail benchmark are imprecise.  A fully interacted system yields 89 supported family-specific edge functionals, but their covariance has rank 50; it is not converted into an invalid 89-degree-of-freedom test or a selected discovery list.  Relaxing Webb availability expands the beta-valid support from 468 to 490 occupations and direct-tail stock from 5.03 to 12.61 percent, yet only five of 23 families then span both tails.  Thin direct overlap is not solely a Webb-support artifact.
\item \textbf{Exact stock accounting separates economic objects.}  Equal-month accounting gives a Q5--Q1 young-stock change of $-0.1098$, an older-stock change of $0.0483$, and hence a young-relative change of $-0.1581$.  An exact log-Shapley accounting assigns $-0.1293$ to within-family ratio changes and $-0.0288$ to changes in older-family weights; audited boundary mass is zero.  These identities describe aggregate stocks.  They neither decompose the nonlinear $-0.1321$ coefficient nor identify causal shares, and no separate young/older regressions are claimed.
\item \textbf{Computer use and family paths are distinct conditioning dimensions.}  On the same 455-occupation support with fixed exposure assignments, the pooled baseline, computer-only, family-month-only, and combined coefficients are $-0.09575$, $-0.19660$, $0.03539$, and $-0.04665$.  The computer-use coefficients are positive in both augmented models, while adding computer use sharply reduces retained target information.  Thus computer use acts as a suppressor in this projection; it does not purify the remaining coefficient as AI.  The repaired pandemic-shortfall exercise uses bounded predictions, a strict 363-occupation population, and linked-household regeneration; its conditioning-movement intervals include zero.  These are correlated diagnostics, not causal horse races.
\item \textbf{Industry, education, cohort, and enrollment checks clarify the target.}  Conditioning occupation--industry microcells on 13 broad-industry post shifts moves $-0.1376$ to $-0.0987$; the paired movement is detected while the residual remains negative.  BA and non-BA coefficients are not statistically distinguishable.  Narrower older comparison bands produce different point estimates, but fixing the older group's exact-age composition changes the pooled coefficient by only 0.0008.  Restricting the young numerator to nonenrolled workers moves it by 0.0082, with an interval containing zero.  Nationally, enrollment among ages 22--25 falls from 0.2336 to 0.2211 while the BA-or-higher share rises from 0.2941 to 0.3121.  These profiles document composition and selection; they do not identify cohort, experience, or enrollment mechanisms.
\item \textbf{Preperiod conclusions depend on both the null and the window.}  Against the short 2022Q4 reference, the full 23-coefficient preperiod null rejects in both structures.  The reference-invariant full-window equality test rejects for the pooled path ($p=.0121$) but not the family path ($p=.0641$).  Excluding 2020Q2--2020Q4 removes the pooled five-percent rejection under both nulls ($p=.0772$ and $p=.0676$); for the family path the short-reference null rejects but within-window equality does not ($p=.0079$ and $p=.1048$).  A linear-slope nondetection cannot validate parallel evolution.  January 2023 remains descriptive chronology, not causal timing.
\item \textbf{Static and dynamic targets are reconciled rather than conflated.}  The static coefficients $S$ are $-0.1321$ pooled and $-0.0217$ conditioned.  The short-reference post functionals $P$ are $-0.1199$ and $-0.2074$, while reference-invariant post-minus-pre functionals $D$ are $-0.1314$ and $-0.0217$.  Paired $D-S$ differences are not detected.  Within the family model, $P-S=-0.1858$ has $p=.074$; what is detected is the change in $P-S$ across conditioning structures ($-0.1980$, $p=.021$), not the corresponding change in $D-S$ ($p=.693$).  HonestDiD applies to $P$, whose conventional intervals include zero, not to $S$ or $D$.
\item \textbf{Endpoint extensions do not identify intensification.}  Relative to July 2026, the pooled movements at the December 2024, September 2025, and December 2025 endpoints are approximately 0.021--0.023 log point and similar in magnitude; the December 2024 paired interval includes zero.  The full window remains canonical and December 2024 is a prominent pre-2025 benchmark, not an outcome-selected replacement.  Family-conditioned endpoint movements remain near zero.
\item \textbf{Flows do not resolve the mechanism.}  Corrected adjacent and twelve-month CPS links give intervals containing zero for employment exit, reported occupation change, and exposure allocation conditional on entry.  Annual young-worker retention is substantially lower than older-worker retention, and observable reweighting does not resolve missing-outcome selection.  Hours and conditional earnings are imprecise.  None of these targets is an employer hiring rate or a mechanical decomposition of the stock coefficient.
\item \textbf{Public benchmarks and ACS sharpen comparability without creating a replication.}  An equal-occupation approximation to BCC's top-two-versus-bottom-three grouping gives $-0.0720$ in CPS and $-0.0168$ after family-month conditioning; their paired movement is not detected at five percent.  Annual ACS reconstructions provide much tighter person-sampling intervals, but family-year results remain sensitive to calendar construction and broader shock uncertainty.  Exact dashboard membership, firm controls, and proprietary worker--firm flows cannot be reproduced.
\item \textbf{The estimation objective is separately audited.}  Giving each occupation total weight one, while retaining final-weight shares within occupation--month cells, changes the pooled coefficient from $-0.1321$ to $-0.0833$.  The paired difference is $0.0488$ ($[-0.0604,0.1581]$), so this design does not detect a difference and does not establish equivalence.
\item \textbf{Architecture and mapping evidence are paired on the current contract.}  Five non-$\lambda$ definitions give pooled native-support coefficients from $-0.0991$ to $-0.0145$, and every alternative-minus-beta paired occupation interval contains zero.  On common AIOE support, some family-conditioned implementation contrasts exclude zero.  A declared bridge-allocation grid leaves the pooled coefficient negative but moves some family-conditioned comparisons across zero.  The paper therefore treats construct choice and mapping implementation as real sensitivities without calling different constructs replications or finite grids identified sets.
\item \textbf{Later reported use is aligned but not identifying.}  A public RPS detailed-occupation workbook retains 121 nonsuppressed occupations.  The unweighted correlation of beta exposure with pooled 2025--26 adoption is 0.504 and the Q5-minus-Q1 adoption difference is 0.267, but Q4 exceeds Q5 and only 26.2 percent of exposure variation remains after major-family means.  The workbook provides neither detailed-cell sampling uncertainty nor adoption timing, so this is descriptive construct validation rather than a mechanism test.
\item \textbf{Mapping matters for population, not just score values.}  A literal SOC-vintage mismatch retains 3.33 percent of computer-and-mathematical employment; documented routing restores 97.7 percent.  Fixed-support score repair changes little, while readmission changes the population.  The invalid merge is an implementation failure, not a robustness specification attributed to another paper.
\item \textbf{Finite numerical targets are certified.}  All 11 declared grouped-binomial specifications pass treatment-rank, recession-direction, independent same-objective solver, fitted-mean, normalized-gradient, fitted-Hessian, and two-sided profile checks at unchanged thresholds.  Boundary nuisance groups in the post-2020 and occupation-seasonality models are profiled on the extended-likelihood face rather than silently deleted.  No audited recession direction moves a reported target, and every focal target is finite.  This resolves numerical existence and separation; it does not certify the sampling model.
\item \textbf{Inference remains conditional.}  Occupation wild scores, 22-family dependence weights, household/sample-unit full refits, elapsed-calendar HAC, and the completed 11-design finite-sample audit answer different stochastic questions.  In the empirical observed-effect design, neither occupation nor family clustering attains nominal coverage for every target; the clean zero-target family-month rejection rates are 6.1 percent, 8.9 percent, and 5.6 percent for occupation weights, family weights, and the cross-fit oracle.  Public Basic Monthly files do not permit complete CPS design-based inference, and these variances cannot be added without a valid decomposition.
\end{enumerate}

\section*{What the paper may and may not claim}

The paper may say that public CPS data contain a sizeable negative young-relative employment-stock association under GPT-4 beta exposure; that its national Q5--Q1 interpretation relies heavily on between-family structure and common restrictions because direct within-family tail support is thin; that broad-family conditioning produces a detected coefficient movement; that exact stock accounting attributes the descriptive aggregate change to both the young numerator and older denominator; that a matched computer-use control strengthens rather than attenuates the coefficient; that preperiod evidence does not validate a causal January-2023 break; and that available flow and adoption evidence does not identify a mechanism.

It may not say that AI caused an individual employment decline, that occupational composition causally explains a percentage of an AI effect, that the residual coefficient is an AI effect, that computerization, reopening, enrollment, or cohort supply is the cause, that nondetection means economic equivalence, that the numerical audit supplies complete CPS-design inference, or that CPS reproduces ADP employer hiring and separation results.

\section*{Current revision state}

The scientific working tree contains the rewritten main manuscript and online appendix; fixed-target stock accounting; support, computerization, cohort, flow, ACS, public-benchmark, mapping, architecture, adoption, dynamic-reconciliation, and inference results; machine-readable registries; module-level code and tests; source and environment manifests; covariance and influence outputs; figure and table generators; failure logs; and numerical-consistency audits.  Restricted microdata, credentials, and direct identifiers are excluded.

The scientific appendix no longer contains the mobility-rematching benchmark or F/G coordinate machinery.  Historical artifacts remain available in the replication archive.  Revision chronology is described honestly as internal git ordering rather than external prospective registration.  Final delivery still requires synchronization of the point-by-point response, diagnosis, source diff, comment matrix, status ledger, and exhibit hashes with the settled manuscript, followed by clean compilation and visual inspection of every PDF.

\section*{Unresolved items}

The following limits remain after feasible execution:

\begin{itemize}[leftmargin=1.8em]
\item observed firm adoption, employer identities, employer-match hiring/separation, and detailed occupation-by-time adoption are unavailable; the RPS extension supplies only a later pooled occupation classification;
\item public files do not supply all confidential survey-design variables or counterfactual age-by-occupation weights for population-control revisions;
\item exhaustive BCC dashboard occupation memberships and tie rules are unavailable, so the public bridge and ACS exercises remain approximations rather than exact reproductions;
\item adjacent occupational coding vintages do not validate age-specific route shares;
\item no external gold-standard label identifies a latent true occupational AI exposure;
\item several direct-tail, flow, subgroup, and conditional estimates remain imprecise;
\item the authorized CPS extracts omit \texttt{EARNWEEK2}, so weekly-earnings coverage after March 2023 cannot be extended without a new extract;
\item affiliation, email, acknowledgments, funding, conflicts, and required journal disclosures require author input and remain \texttt{[AUTHOR TO COMPLETE]}.
\item submission files still require final cross-document consistency checks, bibliography/cross-reference validation, compilation, page-by-page visual quality assurance, refreshed manifests and ledgers, and a clean source diff; these are delivery tasks, not completed scientific analyses.
\end{itemize}

These are not reasons to add unrestricted specification searches or manufacture unavailable inputs.  They define the boundary of the paper's contribution.  Before submission, the author should complete the metadata and disclosure fields and obtain econometric judgment on the maintained shock-cluster interpretation; the present public-data audits do not establish a uniquely correct design-based variance.

\end{document}
